\section{Correctness of NbE}

Next, we want to establish the fact that our \ac{NbE} procedure is correct.
The overall correctness statement consists of two parts:
First, completeness states that any two judgmentally equal terms have identical normal forms.
Second, soundness states that any well-typed term is judgmentally equal to its normal form.
We prove correctness by first defining a semantic analogue to the syntactic judgments, in the form of a \ac{PER} model.
Then each result is proven separately using logical relations.

Since not all \acp{PTS} are normalizing, we need to make additional assumptions about the signature we are working with.
We follow \citet{FP15} in assuming that our \ac{PTS} is predicative.

\begin{definition}[Predicativity]
  A \ac{PTS} signature $(\St, \Ax, \Ru)$ is called \textit{predicative} if there exists a partial order $\leq$ over $\St$ such that:
  \begin{enumerate}
  \item If $(\tmcst s : \tmcst s') \in \Ax$, then $\tmcst s < \tmcst s'$;
  \item If $(\tmcst s_1, \tmcst s_2, \tmcst s_3) \in \Ru$, then $\tmcst s_1 \leq \tmcst s_3$ and $\tmcst s_2 \leq \tmcst s_3$.
  \end{enumerate}
\end{definition}

The main advantage of predicativity is that it restricts our attention to a class of \acp{PTS} that is sufficiently weaker than \rocq itself,
  so that we can mechanize normalization without wrestling with G{\"o}del's incompleteness theorems.
Moreover, predicativity covers important type theories, such as \mltt, so we do not consider it too restrictive.




\subsection{PER model}

We will start by working towards the proof of completeness.
The first things we need to define are the semantic analogues to our syntactic judgments.
This is achieved using a \ac{PER} model, i.e. symmetric and transitive relations that represent types.
Given a \ac{PER} $\per A$, we denote the relation of objects as $\perrel m n {\per A}$, and interpret it as meaning that $m$ and $n$ are equal at type $\per A$.
We write $m \in \per A$ as shorthand for $\perrel m m {\per A}$, and interpret this as meaning that $m$ has type $\per A$.
While \acp{PER} are not reflexive over the entire domain, it follows from symmetry and transitivity that if $\perrel m n {\per A}$, then $m \in \per A$.

The key idea behind defining a \ac{PER} model is that we want to represent the intensional syntactic judgments in an extensional style.
For example, two functions are syntactically equal if their bodies are syntactically equal in an extended context, where the new variable is only equal to itself.
In contrast, two functions are semantically equal if they map any semantically equal inputs to semantically equal outputs.

The main challenge with defining the \ac{PER} model for a \ac{PTS} is providing a \ac{PER} for sorts.
Due to the extensional nature of the judgment, relating two function spaces requires us to extract a \ac{PER} interpretation of their domains and co-domains,
  so that we can meaningfully speak of related inputs and outputs.
This means that a sort essentially needs to be represented as a \ac{PER} of \acp{PER}, while simultaneously being a relation over domain objects.
A common approach to tackle this problem is to define an inductive relation over domain objects simultaneously with a recursive function mapping domain objects to \acp{PER}.
However, this only works for predicative type systems, while our system should be general enough to deal with impredicativity.

The only way to achieve a \ac{PER} model capable of representing impredicativity is to use impredicativity.
We use the approach of \citet{mctt}, where induction-recursion was replaced by impredicativity to allow a mechanization in \rocq.
The idea is to represent sorts (or universes in their setting) as a \ac{PER} on domain objects,
  in such a way that the \ac{PER} representing related objects is produced by the proof that they are related.


\subsubsection{PERs for neutral and normal domain objects}

The main goal of the \ac{PER} model is to relate objects that have the same normal form.
That is, we want to relate objects that readback to the same value when reified at related types.
Since or readback functions are defined for neutral and normal domain objects (rather than just domain objects),
  we cannot directly specify that two domain objects readback to the same normal form.
We do have a readback that works on domain objects, but it is only defined for those objects that represent types.
So, to properly define the \ac{PER} model, we first define relations on neutral and normal domain objects, as well as for types:
\[
\begin{array}{l@{~\iff~}l}
  \perrel e {e'} \perne  & \forall i \in \Nat.\exists E \in \Ne.\ \rbne e E ~ \text{and} ~ \rbne {e'} E \\
  \perrel v {v'} \pernf  & \forall i \in \Nat.\exists V \in \Nf.\ \rbnf v V ~ \text{and} ~ \rbnf {v'} V \\
  \perrel a {a'} \pertyp & \forall i \in \Nat.\exists V \in \Nf.\ \rbtyp a V ~ \text{and} ~ \rbtyp {a'} V
\end{array}
\]
Objects related by any of these \acp{PER} are, by definition (and determinacy of readbacks), guaranteed to readback to the same normal forms.

In addition, we need a \ac{PER} relating reflected neutrals, which we define as follows:
\[
\begin{array}{l@{~\iff~}l}
  \perrel {\dreflect a e} {\dreflect {a'} {e'}} \perneu & \perrel e {e'} \perne
\end{array}
\]
This relation simply checks that the two neutrals are related in $\perne$, and ignores the type annotation of the reflection.
We need $\perneu$ when defining the \acp{PER} for sorts in order to produce a relation that represents types in neutral forms.


\subsubsection{PERs for sorts}

\begin{figure}
  \[
  \begin{array}{c}
    \infer{
      \tmcst s = \tmcst s \in \persort[\tmcst s'] \searrow R
    }{
      (\tmcst s : \tmcst s') \in \Ax
      &
      R \iff \persort
    }
    \qquad
    \infer{
      \dreflect {\tmcst s} e = \dreflect {\tmcst s} {e'} \in \persort \searrow R
    }{
      e = e' \in \perne
      &
      R \iff \perneu
    }
    \\[8pt]
    \infer{
      \dpi a B \rho = \dpi {a'} {B'} {\rho'} \in \persort[\tmcst s_3] \searrow R
    }{
      \begin{array}{c}
        (\tmcst s_1, \tmcst s_2, \tmcst s_3) \in \Ru
        \qquad
        a = a' \in \persort[\tmcst s_1] \searrow R_I
        \qquad
        \mhighlight {R_I \text{ is PER}}
        \\
        \forall \DD :: n = n' \in R_I.\ \evalexpfunc B {\denvext \rho n} = \evalexpfunc {B'} {\denvext {\rho'} {n'}} \in \persort[\tmcst s_2] \searrow R_O(\DD)
        \\
        \forall m, m' \in \Dom.\ m = m' \in R \iff (\forall \DD :: n = n' \in R_I.\ \evalappfunc m n = \evalappfunc {m'} {n'} \in R_O(\DD))
      \end{array}
    }
  \end{array}
  \]
  %
  \caption{Definition of the \acp{PER} for sorts}
  \label{fig:per-sorts}
\end{figure}


Next, we discuss the \acp{PER} representing types and sorts, which are defined simultaneously.
Instead of just a \ac{PER} membership, we define (in \Cref{fig:per-sorts}) a judgment $a = b \in \persort \searrow R$, indicating that $a$ and $b$ are related types of sort $\dcst s$,
  and that $R$ is a \ac{PER} that models both $a$ and $b$.
These relations $R$ are defined using propositional equivalence, which we denote as $R \iff R'$.
This mainly for mechanization purposes, as this style allows us to avoid the use of a propositional extensionality axiom.

Like our syntactic judgments, the definition of $\persort$ relies on the \ac{PTS} specification.
If there is an axiom $(\tmcst s : \tmcst s') \in \Ax$, then $\persort[\tmcst s']$ relates $\tmcst s$ to itself,
  and provides a relation $R$ equivalent to $\persort$ as its \ac{PER} interpretation.
That is, to define $\persort[\tmcst s']$, we need to quantify over \acp{PER} that are equivalent to $\persort$, hence the need for impredicativity.

The most interesting case is the one for function spaces.
Intuitively, given a rule $(\tmcst s_1, \tmcst s_2, \tmcst s_3) \in \Ru$, two function spaces are related in $\persort[\tmcst s_3]$
  if their domains are related in $\persort[\tmcst s_1]$ and their co-domains are related in $\persort[\tmcst s_2]$.
Formally, we can easily express this relation on domains, but not on co-domains, as there might be dependencies to take into account.
To address this, we need to view the output relation, denoted $R_O$ in the rules, as a function from the input relation $R_I$ to \acp{PER}.

To prove correctness of normalization, we need to reason by induction on derivations of $a = b \in \persort \searrow R$.
This requires that the judgment definition is well-founded, which is not the case for impredicative \acp{PTS}.
For example, consider the rule $(\square, \star, \star)$ for polymorphic function spaces in \sysf.
Due to the axiom $\star : \square$, we can construct function spaces of the form $\dpi \star B \rho$.
To verify $\dpi \star B \rho = \dpi \star {B'} {\rho'} \in \persort[\star] \searrow R$, we will have a premise $\star = \star \in \persort[\square] \searrow R_I$,
  where $R_I \iff \persort[\star]$.
But then, to relate the codomains, we need to quantify over objects related by $\persort[\star]$ in order to define $\persort[\star]$, so the definition is not well-founded.

With predicativity, we have a partial order $\leq$ over $\St$ that can be extended to derivations of PER relations.
Namely, we say that $\persort[\tmcst s_1] < \persort[\tmcst s_2]$ provided that $\tmcst s_1 < \tmcst s_2$.
Additionally, we extend the order to $\perneu$ by stating that $\perneu < \persort$ for any $\tmcst s$.
Finally, we need that, if $R_1 \iff R_2$, then $R_1 \leq R$ if and only if $R_2 \leq R$.


\subsubsection{PERs for contexts}

\begin{figure}
  \[
  \begin{array}{c}
    \infer{
      \ctxempty = \ctxempty \in \perctx \searrow R
    }{
      \forall \rho, \rho'.\ \rho = \rho' \in R \iff \rho = \rho' = \denvempty
    }
    \\[8pt]
    \infer{
      \ctxext \Gamma A = \ctxext \Delta B \in \perctx \searrow R
    }{
      \begin{array}{c}
        \Gamma = \Delta \in \perctx \searrow R'
        \qquad
        \mhighlight{R' \text{ is PER}}
        \\
        \forall \DD :: \rho = \rho' \in R'.\ \evalexpfunc A \rho = \evalexpfunc {B} {\rho'} \in \persort \searrow R''(\DD)
        \\
        \forall \rho, \rho', m, m'.\ \denvext \rho m = \denvext {\rho'} {m'} \in R \iff (\exists \DD :: \rho = \rho' \in R' ~ \text{and} ~ m = m' \in R''(\DD))
      \end{array}
    }
  \end{array}
  \]
  %
  \caption{Definition of the \ac{PER} for contexts}
  \label{fig:per-ctx}
\end{figure}


Last, we need to define a \ac{PER} over contexts, denoted $\perctx$ and meant to represent the context equality judgment.
Like the \acp{PER} for sorts, we want membership in $\perctx$ to generate a \ac{PER} representing the related contexts.
These generated \acp{PER} are relations over environments whose domains are the related contexts.
This provides a semantic interpretation of the substitution equality judgment.

We can define $\perctx$ inductively on the shape of contexts.
The rules are shown in \Cref{fig:per-ctx}.
Intuitively, we are simply applying the \acp{PER} for sorts to each type in the context.
The empty context is only related to itself, and the associated \ac{PER} of environments likewise only relates the empty environment to itself.
For context extensions, we relate $\ctxext \Gamma A$ to $\ctxext \Delta B$ provided that $\Gamma$ and $\Delta$ are related,
  and that $A$ and $B$ evaluate to related types in any related environments.
Likewise, the \ac{PER} representing $\ctxext \Gamma A$ relates environments whose heads are related at (the evaluation of) $A$ and tails are related at $\Gamma$.


\subsubsection{Properties of PER model}

The first important property states the our order on relations allows us to do induction on derivations of $a = b \in \persort \searrow R$.

\begin{lemma}
  If $a = b \in \persort \searrow R$, then $R < \persort$.
\end{lemma}
\begin{proof}
  \AG{
    This is problematic: we should be using induction on the derivation, but we can only do this provided that the lemma holds.
    Moreover, we probably need a stronger order relation than just the one on PERs.
    If we consider the rule $(\star, \star, \star)$ for simple function spaces, then the case for $\dpi a B \rho = \dpi {a'} {B'} {\rho'} \in \persort[\star] \searrow R$
      requires induction on the premise $a = a' \in \persort[\star] \searrow R_I$, which can work since $a < \dpi a B \rho$.
  }
\end{proof}

\begin{lemma}
  All the relations defined above really are \acp{PER}, i.e. symmetric and transitive.
\end{lemma}
\begin{proof}
  The cases for $\perne$, $\pernf$, and $\pertyp$ follow from the fact that syntactic equality is an equivalence relation.
  The case for $\perneu$ follows directly from the fact that $\perne$ is a \ac{PER}.
  Once those are established, we can prove by induction that $\persort$ is a \ac{PER} and that if $a = b \in \persort \searrow R$, then $R$ is a \ac{PER}.
  Finally, we prove by induction that $\perctx$ is a \ac{PER} and that if $\Gamma = \Delta \in \perctx \searrow R$, then $R$ is a \ac{PER}.
\end{proof}

The next lemma is essentially a congruence rule for the PERs $\perne$ and $\pernf$.

\begin{lemma}
  If $e = e' \in \perne$ and $v = v' \in \pernf$, then $\dapp e v = \dapp {e'} {v'} \in \perne$.
\end{lemma}
\begin{proof}
  We have $\rbnefunc {\dapp e v} = \tmapp {\rbnefunc e} {\rbnffunc v}$ by definition of readbacks.
  We have $\rbnefunc e = \rbnefunc {e'}$ and $\rbnffunc v = \rbnffunc {v'}$ by definition of $\perne$ and $\pernf$, respectively.
  Then $\rbnefunc {\dapp e v} = \rbnefunc {\dapp {e'} {v'}}$, so $\dapp e v = \dapp {e'} {v'} \in \perne$.
\end{proof}

\AG{We likely need more of these congruence lemmas}


The next lemma states that we only need to know one of the related objects at a sort PER in order to determine what the output relation is, up to relational equivalence.

\begin{lemma}[Relational irrelevance]
  The relations produced from $\persort$ and $\perctx$ are unique up to propositional equivalence.
  Specifically:
  \begin{enumerate}
  \item If $a = b \in \persort \searrow R$ and $a = b' \in \persort \searrow R'$, then $R \iff R'$;
  \item If $a = b \in \persort \searrow R$ and $a' = b \in \persort \searrow R'$, then $R \iff R'$;
  \item If $a = b \in \persort \searrow R$ and $b' = a \in \persort \searrow R'$, then $R \iff R'$;
  \item If $a = b \in \persort \searrow R$ and $b = a' \in \persort \searrow R'$, then $R \iff R'$.
  \end{enumerate}
  And similarly with $\perctx$.
\end{lemma}
\begin{proof}
  By induction on the first derivation.
  In each case, we can apply inversion on the second derivation since rules are syntax-directed
\end{proof}

Next, we have some basic properties for the context PERs.

\begin{lemma}
  If $\Gamma = \Delta \in \perctx \searrow R$, then $|\Gamma| = |\Delta|$.
\end{lemma}
\begin{proof}
  By straightforward induction on the given derivation.
\end{proof}

\begin{lemma}
  \label{lem:initenv}
  Initial environments are related by their context \ac{PER}.
  That is, if $\Gamma = \Delta \in \perctx \searrow R$, then $\denvinit \Gamma = \denvinit \Delta \in R$.
\end{lemma}
\begin{proof}
  By induction on the given derivation.
\end{proof}

Finally, the most important property of the PER model is realizability.
Intuitively, realizability means that all of our PERs sit between $\perne$ and $\pernf$.

\begin{theorem}[Realizability]
  \label{thm:realizability}
  Assume that $a = a' \in \persort \searrow R$.
  Then the followings hold:
  \begin{enumerate}
  \item If $e = e' \in \perne$, then $\dreflect a e = \dreflect {a'} {e'} \in R$;
  \item If $m = m' \in R$, then $\dreify a m = \dreify {a'} {m'} \in \pernf$;
  \item $a = a' \in \pertyp$.
  \end{enumerate}
\end{theorem}
\begin{proof}
  By induction on the derivation $\DD :: a = a' \in \persort \searrow R$.
  This is where we use predicativity of the \ac{PTS}, as this ensures that the induction is well-founded.
\end{proof}



\subsection{Completeness}

\begin{figure}
  \[
  \begin{array}{c}
    \infer{
      \eqctxsem \Gamma {\Gamma'}
    }{
      \Gamma = \Gamma' \in \perctx \searrow R_E
    }
    \qquad
    \infer{
      \wfctxsem \Gamma
    }{
      \eqctxsem \Gamma \Gamma
    }
    \qquad
    \infer{
      \wfexpsem \Gamma M A
    }{
      \eqexpsem \Gamma M M A
    }
    \\[8pt]
    \infer{
      \eqexpsem \Gamma M {M'} A
    }{
      \begin{array}{c}
        \Gamma \in \perctx \searrow R_E
        \\
        \forall \DD :: \rho = \rho' \in R_E.\
        \evalexpfunc A \rho = \evalexpfunc A {\rho'} \in \persort \searrow R
        ~\text{and}~
        \evalexpfunc M \rho = \evalexpfunc {M'} {\rho'} \in R
      \end{array}
    }
    \\[8pt]
    \infer{
      \eqtypsem \Gamma {\tmcst s} {\tmcst s}
    }{
      \text{there is no axiom } (\tmcst s : \tmcst s') \in \Ax
    }
    \qquad
    \infer{
      \eqtypsem \Gamma A {A'}
    }{
      \eqexpsem \Gamma A {A'} {\tmcst s}
    }
    \qquad
    \infer{
      \wftypsem \Gamma A
    }{
      \eqtypsem \Gamma A A
    }
    \\[8pt]
    \infer{
      \eqsubsem \Gamma \sigma {\sigma'} \Delta
    }{
      \begin{array}{c}
        \Gamma \in \perctx \searrow R_\Gamma
        \qquad
        \Delta \in \perctx \searrow R_\Delta
        \\
        \forall \DD :: \rho = \rho' \in R_\Gamma. \
        \evalsubfunc \sigma \rho = \evalsubfunc {\sigma'} {\rho'} \in R_\Delta
      \end{array}
    }
    \qquad
    \infer{
      \wfsubsem \Gamma \sigma \Delta
    }{
      \eqsubsem \Gamma \sigma \sigma \Delta
    }
  \end{array}
  \]
  \caption{Logical relations for completeness}
  \label{fig:logrel-comp}
\end{figure}


The first step to establish completeness is to define logical relations for each judgments,
  which we distinguish from syntactic judgments by using $\semjudg$ instead of $\judg$, e.g. $\wfexpsem \Gamma M A$.
We first define the equality judgments using the PER model, then define the well-formedness by requiring that an object is equal to itself.
Intuitively, we validate equality by evaluating the syntactic objects and checking that they are suitably related in the PER model.
For example, equality between terms is defined as follows:
\[
\infer{
  \eqexpsem \Gamma M {M'} A
}{
  \begin{array}{c}
    \Gamma \in \perctx \searrow R_E
    \\
    \forall \DD :: \rho = \rho' \in R_E.\
    \evalexpfunc A \rho = \evalexpfunc A {\rho'} \in \persort \searrow R
    ~\text{and}~
    \evalexpfunc M \rho = \evalexpfunc {M'} {\rho'} \in R
  \end{array}
}
\]
Remaining judgments are shown in \Cref{fig:logrel-comp}

\begin{theorem}[Fundamental theorem of completeness]
  \label{lem:fundamental-comp}
  If a syntactic judgment holds, then its logical analogue also holds.  Specifically:
  \begin{enumerate}
  \item If $\wfctx \Gamma$, then $\wfctxsem \Gamma$;
  \item If $\eqctx \Gamma \Delta$, then $\eqctxsem \Gamma \Delta$;
  \item If $\wfexp \Gamma M A$, then $\wfexpsem \Gamma M M A$;
  \item If $\eqexp \Gamma M {M'} A$, then $\eqexpsem \Gamma M {M'} A$;
  \item If $\wftyp \Gamma A$, then $\wftypsem \Gamma A$;
  \item If $\eqtyp \Gamma A {A'}$, then $\eqtypsem \Gamma A {A'}$;
  \item If $\wfsub \Gamma \sigma \Delta$, then $\wfsubsem \Gamma \sigma \Delta$;
  \item If $\eqsub \Gamma \sigma {\sigma'} \Delta$, then $\eqsubsem \Gamma \sigma {\sigma'} \Delta$.
  \end{enumerate}
\end{theorem}
\begin{proof}
  By simultaneous induction on the syntactic derivation, using several of the lemmas about the PER model.
\end{proof}

\begin{theorem}[Completeness of \ac{NbE}]
  If $\eqexp \Gamma M {M'} A$, then $\nbeexpfunc \Gamma M A = \nbeexpfunc \Gamma {M'} A$.
\end{theorem}
\begin{proof}
  We have $\wfctx \Gamma$ by \Cref{lemma:pts-presupposition} on $\eqexp \Gamma M {M'} A$, and $\wfctxsem \Gamma$ by \Cref{lem:fundamental-comp}.
  Then $\Gamma \in \perctx \searrow R_E$ by definition of $\wfctxsem \Gamma$, and $\denvinit \Gamma \in R_E$ by \Cref{lem:initenv}.
  We have $\eqexpsem \Gamma M {M'} A$ by \Cref{lem:fundamental-comp}.
  Then $\evalexpfunc A {\denvinit \Gamma} \in \persort \searrow R$ and $\evalexpfunc M {\denvinit \Gamma} = \evalexpfunc {M'} {\denvinit \Gamma} \in R$ by definition of $\eqexpsem \Gamma M {M'} A$.
  Finally, we conclude that $\dreify {\evalexpfunc A {\denvinit \Gamma}} {\evalexpfunc M {\denvinit \Gamma}} = \dreify {\evalexpfunc A {\denvinit \Gamma}} {\evalexpfunc {M'} {\denvinit \Gamma}} \in \pernf$ by \Cref{thm:realizability}.
  Then $\nbeexpfunc \Gamma M A = \rbnffunc[|\Gamma|] {\dreify {\evalexpfunc A {\denvinit \Gamma}} {\evalexpfunc M {\denvinit \Gamma}}} = \rbnffunc[|\Gamma|] {\dreify {\evalexpfunc A {\denvinit \Gamma}} {\evalexpfunc {M'} {\denvinit \Gamma}}} = \nbeexpfunc \Gamma M A$ by definition of $\pernf$.
\end{proof}



\subsection{Kripke gluing model}

We now move our focus to proving soundness of \ac{NbE}.
The proof of soundness follows a similar structure as the one for completeness,
  except that our logical relations will relate a syntactic object to a domain object, instead of two syntactic objects.
In addition, it needs to be stable under weakening.
This is achieved by using a Kripke gluing model in place of the PER model.

The first step in defining the gluing model is to formalize the idea of a weakening substitutions.
We define a new judgment $\wfwk \Delta \sigma \Gamma$, meaning that $\sigma$ is a weakening substitution from $\Gamma$ to $\Delta$:
\[
\begin{array}{c}
  \infer{
    \wfwk \Gamma \sigma \Gamma
  }{
    \eqsub \Gamma \sigma {\subid} \Gamma
  }
  \qquad
  \infer{
    \wfwk \Delta \sigma \Gamma
  }{
    \wfwk \Delta \tau {\ctxext \Gamma A}
    &
    \eqsub \Delta \sigma {\subclo \subwk \tau} \Gamma
  }
\end{array}
\]


We now move on to define the gluing model.
The gluing model is given by a series of relations, each analogous to one of the \acp{PER} defined earlier.
The main relation, denoted $\glusort$, is analogous to our PER $\persort$, but differs from it in two main ways:
First, $\glusort$ is a unary relation on domain types, wheras $\persort$ is binary;
Second, $\glusort$ outputs two relations, wheras $\persort$ only outputs one.
Specifically, a derivation $\glusortrel a P {El}$ outputs a unary predicate $P$ that relates syntactic types to $a$,
  and a binary relation $El$ between syntactic expressions and domain objects of type $a$.

For the PER $\persort$, the output relations where defined simultaneously as $\persort$ itself, with the exception of $\perneu$.
This was necessary to obtain an inductive definition despite the impredicative quantification.
For the gluing relation $\glusort$, we only need to quantify over the PER model, so we will instead define the gluing relations $P$ and $El$ in advance.
We do this using two judgments:
$\glutpjudg \Gamma A P$ states that $A$ is related in $P$,
and $\glutmjudg \Gamma M A m {El}$ states that $M$ and $m$ are related by $El$.

We start by defining two predicates classifying neutral types and terms.
These are denoted $\gluneutp a$ and $\gluneutm a$, respectively, and they are analogous to our \ac{PER} $\perneu$.
\[
\begin{array}{c}
  \infer{
    \glutpjudg \Gamma A {\gluneutp a}
  }{
    \wfexp \Gamma A {\tmcst s}
    &
    \forall \Delta \in \Ctx. \forall \sigma \in \Sub. \wfwk \Delta \sigma \Gamma \implies \eqexp \Delta {\tmclo \sigma A} {\rbnefunc[|\Delta|] a} {\tmcst s}
  }
  \\[8pt]
  \infer{
    \glutmjudg \Gamma M A {\dreflect b e} {\gluneutm a}
  }{
    \begin{array}{c}
      \glutpjudg \Gamma A {\gluneutp a}
      \qquad
      \wfexp \Gamma M A
      \qquad
      e \in \perne
      \\
      \forall \Delta \in \Ctx. \forall \sigma \in \Sub. \wfwk \Delta \sigma \Gamma \implies \eqexp \Delta {\tmclo \sigma M} {\rbnefunc[|\Delta|] e} A
    \end{array}
  }
\end{array}
\]
Note that, for $\gluneutp a$, we use the neutral readback and not the type readback, which ensures that $a$ is indeed a neutral type.
Note also that, for $\gluneutm a$, we reflect the neutral domain object at some type $b$, not necessarily $a$.
This is similar to the \ac{PER} $\perneu$, where type annotations are ignored by the relation.

Next, we define the relations for function spaces.


\subsection{Soundness}






\begin{theorem}[Soundness of \ac{NbE}]
  If $\wfexp \Gamma M A$, then $\eqexp \Gamma M {\nbeexpfunc \Gamma M A} A$.
\end{theorem}
